民爆不是锂价周期的反向对冲工具，但它提供了雅化第二个、盈利波动较低的利润池。投资判断的关键不在于“公司有民爆牌照”，而在于产品、矿服和运输的毛利结构是否能在锂盐回落时维持现金贡献。我们只给予已审计的历史收入和毛利信用；未披露的 2026 年项目订单、回款和新增许可利用率不进入基准估值。

\section{2025 年利润池：产品毛利高，矿服回款仍需单独验证}

民爆产品 2025 年收入 21.41 亿元，毛利 10.11 亿元，毛利率 47.22\%；爆破服务收入 11.75 亿元，毛利 2.31 亿元，毛利率 19.67\%。两者合计收入 33.16 亿元、毛利 12.42 亿元，运输收入另为 4.03 亿元。产品与服务的毛利差意味着“民爆稳定”不能被简化为一个合并倍数：服务业务的项目进度和回款质量对现金含量更关键。

\begin{exhibitbox}[表 4-1：民爆与矿服的历史利润池]
\begin{tabularx}{\exhibitboxwidth}{L{3.4cm}R{2.3cm}R{2.3cm}R{2.2cm}X}
\toprule
2025A 业务 & 收入 & 成本 & 毛利率 & 基准情景处理 \\
\midrule
民爆产品 & 21.41 亿元 & 11.30 亿元 & 47.22\% & 给予历史经营底座信用，不预设许可扩张 \\
爆破及矿服 & 11.75 亿元 & 9.44 亿元 & 19.67\% & 给予历史规模信用，项目回款不作增量假设 \\
合计 & 33.16 亿元 & 20.74 亿元 & 37.46\% & 作为锂盐熊市的防御性利润池 \\
危化运输 & 4.03 亿元 & 未单独列示 & 未单独列示 & 不单独赋予估值溢价 \\
\bottomrule
\end{tabularx}
\end{exhibitbox}
\sourcenote{公司《2025 年年度报告》分产品收入、成本披露；合计和毛利率为 AStock 复算。}

该利润池的投资含义是下行保护而非上行期权。基准情景不使用新项目、并购、海外矿服或许可证扩大来提高 2026 年民爆利润；这避免把尚未披露的订单当作对锂盐利润的补偿。

\section{许可、区域与一体化服务：解释壁垒，不等于新增收入}

公司披露截至 2026 年 4 月拥有炸药许可产能 33 万吨、雷管许可产能近 9,000 万发；民爆产品执行以销定产，爆破服务覆盖矿山和基建项目。许可和区域网络构成进入壁垒，且一体化服务可提高客户黏性，但许可产能不是当期销量，5--10 年战略合作也不是已确认订单。

\begin{exhibitbox}[表 4-2：民爆防御性信用的准入规则]
\begin{tabularx}{\exhibitboxwidth}{L{3.5cm}X X}
\toprule
项目 & 可以进入基准模型的内容 & 只能作为监控或牛情景的内容 \\
\midrule
产品业务 & 历史收入、成本、毛利和季节性 & 33 万吨许可所对应的新增销量与利润 \\
爆破及矿服 & 历史收入、毛利与业务模式 & 新项目金额、剩余 backlog、实际回款与毛利改善 \\
客户与区域 & 已披露的行业/区域覆盖与合作框架 & 客户份额、订单价格、合同剩余期限和收入确认 \\
运输业务 & 历史收入存在与业务协同性 & 独立利润、现金流和估值倍数 \\
\bottomrule
\end{tabularx}
\end{exhibitbox}

\section{下一次验证：利润托底必须同时通过现金检验}

民爆产品通常受许可、区域物流和下游矿山/基建开工约束，第一季度还受春节和天气影响。正式 H1 应同时核验分部收入、毛利、应收和经营现金流。若民爆利润增长伴随更高应收或项目占款，它只能稳定会计利润，不能稳定自由现金流；这会降低其对合并估值倍数的支撑。
